\section{Multi-Agent Systems：并行化、专业化与通信协议}

单 agent 面对长任务时会受上下文、工具数量和错误累积限制。多 agent 系统尝试把任务拆分给不同 worker，由 manager 协调结果。它可以并行化和专业化，也引入新的分布式系统问题：任务划分、消息一致性、重复执行、同步等待、权限传播和故障恢复。

\subsection{为什么使用多个 Agents}

第一张图用多个机器人表示自主系统；完整 build 列出 parallelization 和 task specialization。并行化只有在子任务依赖较弱时才缩短时延；若每个 worker 都等待同一共享状态，通信开销会抵消收益。

这一节从单 agent 的 memory 与 personalization 转向协作边界。拆分前应先画依赖图：检索多个独立来源可以并行，写最终报告必须等待证据汇总；不同 worker 可以各自使用专门工具，但共享用户身份和付款权限则需要中央控制。读图时不要把机器人数量当作能力指标，而要问每个 agent 是否承担独立职责、输出能否合并、冲突由谁裁决，以及一个 worker 超时或返回恶意内容时 manager 是否会隔离它。

\lecturefigure{slide-52-multi-agent-autonomous-systems.jpg}{多 Agent 自主系统的基本形态}{00:39:49}

\lecturefigure{slide-53-why-multiagent-systems.jpg}{多 Agent 的两项主要动机：并行化与专业化}{00:40:48}

若任务被分为 $K$ 个子任务，串行时延为 $\sum_k T_k$；理想并行时延近似为
\[
T_{\text{parallel}}\approx \max_k T_k+T_{\text{coord}}+T_{\text{merge}}.
\]
$T_{\text{coord}}$ 是调度与通信，$T_{\text{merge}}$ 是结果合并与验证。若协调成本太高或最长子任务占主导，多 agent 不会自动更快。

\begin{warningbox}{“多 Agent”不是把同一 Prompt 复制多份}
没有角色边界、共享状态协议和冲突解决的多 agent，通常只是提高 token 成本。每个 worker 应有明确输入、输出 schema、权限、截止时间和失败处理。
\end{warningbox}

\subsection{Hierarchy 与 syncing primitives}

层级图把用户请求交给 manager，再分派给多个 worker。Manager 的责任不只是转发消息，还要规划依赖、控制预算、判断结果冲突，并在 worker 失败时重新分配。所谓 \term{syncing primitives} 是让多个 agent 对共享状态达成可操作一致的机制，例如锁、版本号、barrier、消息确认和幂等任务 ID。

\lecturefigure{slide-54-agent-hierarchy.jpg}{Manager--worker 层级与 Agent 间信息流}{00:42:21}

\begin{lstlisting}[caption={Manager--worker 任务合同}]
task = {
  "task_id": stable_id,
  "goal": scoped_goal,
  "inputs": immutable_inputs,
  "permissions": allowed_tools,
  "depends_on": prerequisite_task_ids,
  "deadline": deadline,
  "output_schema": schema,
  "retry_policy": retry_policy
}
\end{lstlisting}

\subsection{MCP 与 A2A：连接标准不等于安全保证}

协议 slide 把 MCP（Model Context Protocol）用于模型、工具和数据连接，把 A2A（Agent-to-Agent）用于 agent 通信。MCP 的价值是统一发现和调用接口，减少每个应用为每个工具写定制 glue code；A2A 则需要定义身份、消息、任务状态与协作语义。二者解决互操作性，不自动解决授权和正确性。

\lecturefigure{slide-55-mcp-a2a-protocols.jpg}{MCP 连接模型与工具，A2A 支持 Agent 间协作}{00:43:44}

\begin{knowledgebox}{首次使用：Protocol、Schema 与 Authorization}
\textbf{Protocol} 定义参与方怎样通信；\textbf{schema} 定义消息字段及类型；\textbf{authorization} 决定谁可执行什么动作。MCP server 暴露工具，不代表任何 client 都应拥有调用权限；权限仍需最小化、可审计且可撤销。
\end{knowledgebox}

\subsection{为什么 MCP 被比作 USB-C}

第一张 MCP 图强调统一连接模型、数据与工具；最终 build 进一步列出 resilience、dynamic tool discovery 和 intent-based execution。读图时应把“吸收 API 变化”理解为抽象层隔离，而不是接口永不出错：server schema 变化、语义变化和权限变化仍需版本管理和兼容测试。

\lecturefigure{slide-56-mcp-usbc.jpg}{MCP 的 USB-C 类比：统一模型与工具连接}{00:45:11}

\lecturefigure{slide-57-mcp-api-benefits.jpg}{MCP 相对点对点 API 集成的韧性与动态发现}{00:45:46}

协议标准化还会改变故障定位。点对点集成失败时，应用、模型和 API 的责任常混在一起；引入协议层后，可以分别检查 server 是否可发现、schema 是否匹配、权限是否足够、调用是否超时、返回是否通过验证。代价是协议层本身也成为关键依赖，必须版本化、监控并防止恶意 server 注入误导性工具描述。所谓 resilience 来自清晰分层和测试，不是来自一个统一名字。

\teachervoice{00:38:00--00:46:09，讲者把多 agent 的价值归纳为 parallelization 与 specialization，并强调 hierarchy、synchronization、MCP、A2A、resilient integrations 和 dynamic discovery。课堂提示：自然语言本身含糊，通信必须有结构化协议和状态。}

\begin{importantbox}{协议层的最小安全清单}
工具发现应与权限授予分离；敏感工具要求显式 scope；每次调用包含调用者身份和任务 ID；输出需要 schema 验证；协议升级要做兼容测试；高风险动作应要求用户确认或双重控制。
\end{importantbox}

\subsection{本章小结}

多 agent 系统把单模型问题转成分布式系统问题。并行化和专业化只有在任务可拆、通信可控时有收益；manager--worker 合同、同步 primitives、MCP/A2A schema、身份与权限共同决定系统是否可靠。

